\begin{theorem}[Theorem 4]
Let $n\ge2$, $p>1$, $0<\lambda\le\Lambda$, $a>0$, $0<\nu_0<\nu_1$ and
$\delta>0$, suppose the Sobolev--Poincar\'e inequality on balls holds in dimension $n$
(\noderef{CitedSobolevPoincareBall}), and assume the cited regularity result
(C1) (\noderef{CitedRegularityEpsPos}), the cited solvability result (C2)
(\noderef{CitedDirichletW1p}) and the cited Dirichlet result (C3) with
H\"older boundary data (\noderef{CitedDirichletHolder}) in the universal forms
described below, and assume the cited Campanato embedding (C5) at the exponent
$n/(n-1)$ on every ball (\noderef{CitedCampanatoEmbedding}), in the universal
form described below.  Then there is a number $\lambda_0$ with
$0<\lambda_0\le p$, depending only on these quantities, such that for all
$c>0$, $d>0$, $M\ge0$, $W\ge0$ and every modulus $\sigma\colon\R\to\R$ there is
$N\ge0$ with the following property.

Let $0\le\varepsilon\le1$, let $\Omega\subseteq\R^n$ be open, let $\Omega'$ be a
open, nonempty and bounded and a bounded Lipschitz domain
(\noderef{IsBoundedLipschitzDomain}), with $\overline{\Omega'}\subseteq\Omega$ and
$d\le d(\Omega',\partial\Omega)$ --- ``$\Omega'$ is a bounded Lipschitz domain'' being the hypothesis the paper states at l.~1274, now carried as the predicate \noderef{IsBoundedLipschitzDomain} and no longer only as the three clauses open, nonempty and bounded; those three are kept beside it, being implied by it, so that the proof sites discharging them by name read unchanged ---
and assume the
cited Campanato embedding (C5) on $\Omega'$ at the exponent $p$
(\noderef{CitedCampanatoEmbedding}) --- and, at the second exponent
$n/(n-1)$, the instance of (C5) on $\Omega'$ that the Morrey theorem this proof
consumes requires in its turn (\noderef{CitedCampanatoEmbedding} again): both
instances are hypotheses here, bound with $\Omega'$ because (C5) mentions
$\Omega'$ and so cannot be hoisted above it, and assume the cited H\"older
estimates (C4) (\noderef{CitedHolderOscillationMinimum}).  Let $f$ satisfy
(F1) of the thesis--(F3) of the thesis on $\Omega$ with parameters
$p,\varepsilon,\lambda,\Lambda,a,\nu_0,\nu_1$ and the modulus $\sigma$ fixed
above (\noderef{VariationalStructureF123}), and in addition (F4) of the thesis:
$\sigma(t)\le c\,t^{\delta}$ for $t\ge0$.  Let
$u\in W^{1,p}(\Omega)\cap L^\infty(\Omega)$ with weak gradient $g=\nabla u$
(\noderef{MemW1pOn}), $|u|\le M$ a.e.\ and
$\int_\Omega|u|^p\,dx+\int_\Omega|\nabla u|^p\,dx\le W^p$ (that is,
$\|u\|_{W^{1,p}(\Omega)}\le W$), be a local minimum for $F$ in
$\Omega$ (\noderef{IsLocalMinimumOn}).  Then
\begin{itemize}
\item[(i)] $\|D_iu\|_{\mathcal L^{p,\,n+\lambda_0}(\Omega')}\le N$ for every
  $i=1,\dots,n$, i.e.\ $\nabla u\in\mathcal L^{p,\,n+\lambda_0}(\Omega')$
  componentwise (\noderef{CampanatoBoundOn}); and
\item[(ii)] $u$ agrees almost everywhere on $\Omega'$, together with its
  gradient, with a function $w$ that is differentiable at every point of
  $\Omega'$ with $\nabla w=g_w$ there and whose gradient $g_w$ is continuous and
  $\beta$-H\"older on $\overline{\Omega'}$, $\beta=\lambda_0/p$
  (\noderef{MemC1HolderOn} at the pair of sets $(\Omega',\overline{\Omega'})$);
  i.e.\ $u\in C^{1+\beta}(\Omega')$ for some $\beta$ with $0<\beta\le1$, the
  H\"older seminorm of the gradient being taken on $\overline{\Omega'}$ as the
  Note's Definitions paragraph at l.~278 prescribes.
\end{itemize}
The quantifier order records the dependences asserted by the paper: $\lambda_0$
and $\beta=\lambda_0/p$ are chosen as functions of
$\lambda,\Lambda,\nu_0,\nu_1,n,p,a$ and $\delta$ alone, before the constant $c$
of (F4) of the thesis, $d(\Omega',\partial\Omega)$, $\|u\|_{L^\infty(\Omega)}$,
$\|u\|_{W^{1,p}(\Omega)}$ and the modulus $\sigma$ of (F3) of the thesis; the norm
$N$ is chosen after all five of those in addition.  The modulus $\sigma$ is
bound with $c$, $d$, $M$ and $W$ rather than with the remaining data because the
norm genuinely depends on it: $N$ is built from the Morrey constant of
Theorem~Theorem~3 of the thesis, and the paper's own remark on that constant, at l.~1283,
lists ``the function $\sigma$ given in (F3) of the thesis'' among its dependences.
Theorem~Theorem~4 of the thesis says its norm depends ``in addition to the previous
constants'' on $d(\Omega',\partial\Omega)$, $\|u\|_{L^\infty}$ and
$\|u\|_{W^{1,p}}$, and ``the previous constants'' are those of
Theorem~Theorem~3 of the thesis, $\sigma$ among them; so naming $\sigma$ here reports the
paper's dependence rather than adding one.  Nothing is lost on the side that matters:
$\lambda_0$ and $\beta$ are still bound \emph{before} $\sigma$, which is what
the paper's own list for them at l.~1353--1355 asserts, and the paper fixes $f$,
hence $\sigma$, before Theorem~Theorem~4 of the thesis, which is stated ``in the situation
of Theorem~Theorem~3 of the thesis''.  In particular the H\"older \emph{exponent} $\delta$ of (F4) of the thesis
may influence $\lambda_0$ and $\beta$, while its \emph{constant} $c$ may not:
this is what the paper asserts, and the proof of Theorem~Theorem~4 of the thesis confirms it
by producing a $\delta'>0$ ``depending only on $\delta$, $\lambda$, $\Lambda$,
$\nu_0$, $\nu_1$, $n$, $p$ and $a$'' with no $c$.  As the Note for the
formalization prescribes, the vector field of part (i) is $\nabla u$ taken
componentwise.
\noindent\textbf{On the $W^{1,p}$ datum.}  The hypothesis
$\int_\Omega|u|^p\,dx+\int_\Omega|\nabla u|^p\,dx\le W^p$ is a bound on the
genuine Sobolev norm $\|u\|_{W^{1,p}(\Omega)}$, which is the quantity the
paper's dependency list names.  It must not be replaced by a bound on
$\int_\Omega(1+|\nabla u|^2)^{p/2}dx$: that integrand is pointwise at least $1$,
so such a bound would force $W\ge|\Omega|$, would be unsatisfiable whenever
$|\Omega|=\infty$ --- and the paper states this result for an arbitrary open
$\Omega\subseteq\R^n$, so $\Omega=\R^n$ is an instance --- and would let the
conclusion's norm degrade with $|\Omega|$, which is not on the paper's list.
The weighted integral is still available locally where the proof needs it: since
$(1+t^2)^{p/2}\le2^{p/2}(1+t^p)$ for $t\ge0$,
\[
\int_{B_R}(1+|\nabla u|^2)^{\frac p2}dx
\;\le\;2^{\frac p2}\Bigl(|B_R|+\int_{B_R}|\nabla u|^p\,dx\Bigr)
\;\le\;2^{\frac p2}\bigl(c_nR^{\,n}+W^p\bigr),
\]
in which every term is controlled by $n$, $p$, $R$ and $W$ alone.

The hypothesis \noderef{CitedSobolevPoincareBall} is the cited result (C6) of the
paper's Note for the formalization, carried here as an explicit argument because
the proof uses it --- it enters through the constants of \noderef{MinimumMorreyRegularity} and \noderef{GradientOscillationEstimate} --- which is what the
Note permits for the results the thesis quotes without proof.  Although it is
bound before the existential for $\lambda_0$, it introduces no dependence of
$\lambda_0$ on any quantity outside the list the paper names: (C6) is a statement
about $n$ alone, and $n$ is already bound first, so the hypothesis carries no
numerical content of its own.

\medskip
\noindent\textbf{The cited Campanato embedding on balls.}  The cited result
(C5) is assumed here at the exponent $n/(n-1)$ on \emph{every} ball
$B_\rho(y)\subseteq\R^n$ with $\rho>0$ (\noderef{CitedCampanatoEmbedding}), in
addition to the two per-instance assumptions on $\Omega'$ above.  A ball is a
bounded Lipschitz domain, so this is an instance of (C5) exactly as the paper's
Note states it, universally quantified over centre and radius in the same way
\noderef{CitedSobolevPoincareBall} quantifies (C6) --- and the Note allows the
cited results to ``be taken as explicit hypotheses \ldots\ where they are
used''.  It is used in exactly one place: the proof needs the H\"older
continuity of $u$ at a \emph{fixed} exponent, with a constant uniform over the
centres $x_0\in\Omega'$, on balls \emph{around} the points of $\Omega'$ rather
than on $\Omega'$ itself, and it gets it by applying part~(iii) of
Theorem~Theorem~3 of the thesis (\noderef{MinimumMorreyRegularity}) on finitely many balls
covering $\overline{\Omega'}$, each of which needs its own instance of (C5).
This cannot be replaced by the exponent of (C4): the Note allows (C4)'s
$\alpha$ to depend on $M=\|u\|_{L^\infty}$, while $\lambda_0$ is bound before
$M$.  Like (C6), the hypothesis is bound before the existential for $\lambda_0$
and introduces no dependence of $\lambda_0$ on anything outside the paper's
list: it mentions only $n$, already bound, and is universally quantified over
the ball.

\medskip
\noindent\textbf{The three cited results for the frozen operator.}  The cited
regularity result (C1) (\noderef{CitedRegularityEpsPos}), the cited solvability
result (C2) (\noderef{CitedDirichletW1p}) and the cited Dirichlet result (C3)
with H\"older boundary data (\noderef{CitedDirichletHolder}) are assumed here
for \emph{every} field $A'\colon\R^n\to\R^n$ and all parameters
$\varepsilon',\gamma_0',\gamma_1'$ with $\gamma_0'>0$ and $\gamma_1'>0$ --- and,
for (C1), for every open $\Omega_0\subseteq\R^n$ as well --- which are the forms
these three results have in the paper's Note for the formalization, where they
``may be taken as explicit hypotheses \ldots\ where they are used''.

The universal shape is forced rather than convenient, and it is forced by the
same frozen operator as in Theorem~Theorem~3 of the thesis: the proof reaches (C1) and (C2)
through \noderef{GradientSupEstimate} (Theorem~1) and all three through
\noderef{MinimumMorreyRegularity} (Theorem~3), each of which carries them in
exactly this universal form and binds them before its own existential, because
the operator to which they are applied is frozen at a centre and a value chosen
inside the proof.  Like (C6), none of the three introduces a dependence of
$\lambda_0$ on any quantity outside the list the paper names: each mentions only
$n$ and $p$, both already bound, and each is universally quantified over its
remaining parameters, so none carries numerical content of its own.  The guards
$\gamma_0'>0$, $\gamma_1'>0$ are the thesis' own standing positivity conditions
on the structure constants, and the openness guard on (C1) is the domain
condition \noderef{CitedRegularityEpsPos} is stated for; (C2) and (C3) need no
such guard because \noderef{CitedDirichletW1p} and
\noderef{CitedDirichletHolder} are statements about balls and fix their own
domains.
\end{theorem}
\begin{proof}
The route is clause (B) of \noderef{ComparisonSchemeCampanatoCore}, which was
built for this theorem and Theorem~Theorem~6 of the thesis and which carries the
construction of the comparison map and the oscillation estimate (4.16) of the thesis
itself; the present proof supplies only what is proper to Theorem~Theorem~4 of the thesis.

\emph{Step 1: the exponents $\lambda_0$ and $\beta$, chosen before $c$, $d$, $M$,
$W$ and $\sigma$.}
\begin{enumerate}
\item[(1a)] Put $r=\nu_1/\nu_0>0$.
\item[(1b)] Apply \noderef{ComparisonGradientEstimateMin} (Lemma~4.1) with the
  parameters $n,p,\lambda,\Lambda,a,\nu_0,\nu_1$ --- its hypotheses $2\le n$,
  $1<p$, $0<\lambda\le\Lambda$, $0<a$, $0<\nu_0<\nu_1$, the
  Sobolev--Poincar\'e inequality (\noderef{CitedSobolevPoincareBall}) and the
  universal cited result (C3) (\noderef{CitedDirichletHolder}) are hypotheses of
  the present theorem --- and take clause (B) of its conclusion, read at the
  H\"older exponent $\alpha=\tfrac12$ and at the present $\delta$: its guards
  $0<\tfrac12\le1$ and $0<\delta$ hold.  It produces $\delta'>0$ depending only
  on $n,p,\lambda,\Lambda,a,\nu_0,\nu_1,\alpha$ and $\delta$ --- every one of
  them on the list the present $\lambda_0$ is allowed to depend on.  The
  exponent $\alpha=\tfrac12$ is a fixed number and not (C4)'s exponent; this is
  the whole reason $\delta'$, and with it $\lambda_0$, is free of
  $M=\|u\|_{L^\infty}$, and Step~(2a$'$) is what pays for it.

  \emph{That node has two clauses and only clause \textup{(B)} is used here.}
  Its clause (A) is Lemma~4.1 as the thesis prints it: it produces a modulus
  $\sigma_2$ and a radius $R_1$ bound \emph{after} $M$, $d$ and $\sigma$, and no
  $\alpha$ or $\delta$ appears anywhere in its scope, because Lemma~4.1 does not
  assume (F4) of the thesis.  So clause (A) carries no decay rate for $\sigma_2$, and
  a rate is exactly what Theorem~Theorem~4 of the thesis needs (pm-0043).  Clause (B) is the
  clause with the $\alpha,\delta$ prefix and the rate $\delta'$; nothing below
  reads clause (A), and in particular the $\sigma_2$ and $R_1$ it produces are
  not the $R_1$ of (2d), which is built from clause (B)'s $R_1'$.
\item[(1c)] Put
  \[
  \tau^\sharp=\frac12\cdot\frac{p\alpha\delta'}{n+p\alpha}>0,
  \qquad\text{so}\qquad
  \tau^\sharp<\frac{p\alpha\delta'}{n+p\alpha},
  \]
  a function of $n,p,\alpha,\delta'$ alone, and apply
  \noderef{ComparisonSchemeCampanatoCore} at the head parameters
  $n,p,\alpha=\tfrac12,\delta'$, $\tau=\tau^\sharp$ and
  \[
  c_1=1,\qquad c_3=1 .
  \]
  Its head guards $2\le n$, $1<p$, $0<\alpha\le1$, $0<\delta'$, $0<c_1$,
  $0<c_3$, $0<\tau^\sharp$ and $\tau^\sharp<p\alpha\delta'/(n+p\alpha)$ all
  hold.  \emph{The values $c_1=c_3=1$ are deliberate and they are what keeps the
  present $\lambda_0$ clean.}  Those two head parameters are read only by clause
  (A) of that node --- $c_1$ in clause (A)'s Morrey hypothesis and $c_3$ in its
  (4.16) of the thesis hypothesis --- and clause (B), which is the clause used here,
  reads neither: it carries its own Morrey constant $c_2$ and builds its own
  (4.16) of the thesis.  Passing the \emph{actual} Morrey constant of Theorem~3 as
  the head's $c_1$ would be legal but fatal, because that constant carries
  $d(\Omega',\partial\Omega)$, $\|u\|_{L^\infty}$, $\|u\|_{W^{1,p}}$ and
  $\sigma$, and the node's $\lambda_2$ sits inside the head prefix and so would
  inherit all four, whereas the present $\lambda_0$ is bound before $d$, $M$,
  $W$ and $\sigma$.  Any fixed positive numbers would do; $1$ is the simplest.
\item[(1d)] Take clause (B) at the ratio $r$ of (1a) and at the H\"older
  exponent $\alpha_H=\tfrac12$.  That clause quantifies over $r$ and
  $\alpha_H$ \emph{universally}, so the value $\tfrac12$ is our choice and not
  one the node prescribes; what the node asks of the choice is the three guards
  $0<r$, $0<\alpha_H$ and $\alpha_H\le1$, and they hold --- $r=\nu_1/\nu_0>0$ by
  (1a) since $0<\nu_0<\nu_1$, and $0<\tfrac12\le1$.  The value $\tfrac12$ must
  agree with the exponent at which (2a$'$) will supply the H\"older datum and
  with the $\alpha$ of (1b), and it does.  Supply as well the
  Sobolev--Poincar\'e inequality
  (\noderef{CitedSobolevPoincareBall}) and the three universal cited results
  (C2) (\noderef{CitedDirichletW1p}), (C1) (\noderef{CitedRegularityEpsPos}) and
  (C3) (\noderef{CitedDirichletHolder}), all four of which are hypotheses of the
  present theorem in exactly the forms that clause requires --- (C1) alone also
  quantifying over every open set.  It produces
  $\alpha_2\in(0,1]$.  Put $\delta_2=\delta'$,
  $\theta_2=\delta_2/(n+p\alpha_2)$ and
  \[
  \tau_2=\min\Bigl\{\tfrac12p\alpha_2\theta_2,\ \tfrac n2,\ \tfrac p2\Bigr\}.
  \]
  Clause (B) binds $\delta_2$ and $\tau_2$ universally behind three guards, and
  all three hold: $0<\delta_2$, because $\delta_2=\delta'>0$ by (1b);
  $0<\tau_2$, because $\alpha_2>0$, $\theta_2=\delta_2/(n+p\alpha_2)>0$ and
  $p>1$ make $\tfrac12p\alpha_2\theta_2>0$, while $n\ge2$ makes $n/2>0$ and
  $p>1$ makes $p/2>0$; and
  $\tau_2<p\alpha_2\delta_2/(n+p\alpha_2)=p\alpha_2\theta_2$, because
  $\tau_2\le\tfrac12p\alpha_2\theta_2<p\alpha_2\theta_2$.  \emph{The third term
  of the minimum is there for} (2a$'$), \emph{and it costs nothing here}: the
  guard clause (B) imposes is the strict inequality just checked, which only
  gets easier as $\tau_2$ shrinks.  Note also
  $\tau_2\le n/2<n$, which is not a guard of this clause but the guard
  Theorem~3 imposes on its own exponent in (2a$'$) and (2c) below; and
  $\tau_2\le p/2<p$ together with $\tau_2>0$, which is what (2a$'$) needs in
  order to meet (C5)'s admissibility and to make the exponent $1-\tau_2/p$ it
  delivers at least $\tfrac12$.  The clause
  then produces $\lambda_2\in(0,p)$ and $R_3\in(0,1]$.
\item[(1e)] Put $\lambda_0=\lambda_2$ and $\beta=\lambda_0/p$.  Then
  $0<\lambda_0$ and $\lambda_0\le p$, because $\lambda_2<p$; and
  $0<\beta<1$, so in particular $0<\beta\le1$ as part (ii) asserts.  (This is
  where the strict bound $\lambda_2<p$ is used: at $\lambda_2=p$ the exponent
  $\lambda=n+\lambda_0$ would fail (C5)'s requirement $\lambda<n+p$ in Step 4.)
  The dependence is right: $\lambda_2$ depends only on
  $n,p,r,\alpha_H,\delta_2,\tau_2$ and on the head's
  $\alpha,\delta',\tau^\sharp,c_1,c_3$, and each of those is a function of
  $n,p,\lambda,\Lambda,a,\nu_0,\nu_1$ and $\delta$ alone --- $r=\nu_1/\nu_0$,
  $\alpha=\alpha_H=\tfrac12$, $\delta'=\delta_2$ by (1b), $\tau^\sharp$ and
  $\tau_2$ by their formulas, and $c_1=c_3=1$.
\end{enumerate}

\emph{Step 2: the case split on (F4) of the thesis, the constants, and the four
per-centre hypotheses clause \textup{(B)} asks for.}  Clause (B) binds sixteen
hypotheses in all, enumerated in Step~3; twelve of them are hypotheses of the
present theorem or immediate from them, and the four that need work are the ones
quantified over the centre $x_0\in\Omega'$ --- the frozen bundle, the H\"older
datum, the Morrey bound and the smallness.  Those four are (2b), (2a$'$), (2c)
and (2d).

Fix $c,d,M,W$ and $\sigma$.  \emph{The data
$\varepsilon,\Omega,\Omega',f,u,g$ are not fixed yet, and must not be}: the
statement binds $N$ above them, so every constant produced between here and
Step~4 --- $N$ included --- has to be a function of
$n,p,\lambda,\Lambda,a,\nu_0,\nu_1,\delta$, the head's cited hypotheses, and
$c,d,M,W,\sigma$ alone.

\emph{The case split on (F4) of the thesis.}  The condition (F4) of the thesis, namely
$\sigma(t)\le c\,t^{\delta}$ for every $t\ge0$, is a statement about $\sigma$,
$c$ and $\delta$ alone --- it mentions no datum --- and all three are bound
above the existential for $N$: $\delta$ at the head of the theorem, with
$0<\delta$, and $c$ and $\sigma$ in the block just fixed.  So we may decide it
here and argue by cases.

If (F4) of the thesis \emph{fails}, take $N=0$.  The only demand the statement makes
on $N$ outside the inner implication is $N\ge0$, and $0\ge0$.  The inner
implication is then satisfied vacuously, because (F4) of the thesis is itself one of
its hypotheses: the present theorem assumes, among the hypotheses on the data,
that $\sigma(t)\le c\,t^{\delta}$ for $t\ge0$.  So for any
$\varepsilon,\Omega,\Omega',f,u,g$ satisfying all of those hypotheses
(F4) of the thesis would hold, contradicting the branch; hence no such data exist and
there is nothing to prove.  Nothing further is needed in this branch, and no
dependence list is affected, $0$ depending on nothing at all.

The rest of the proof runs inside the branch in which (F4) of the thesis
\emph{holds}.  It is fixed once, here, and every constant below is produced
under it.  It contributes no new quantity to any dependence list, being a
property of the already-bound $\sigma$, $c$ and $\delta$.

\emph{The order in which the constants are produced, and where each application
splits.}  In this branch the constants are built in the following order, each
from the head parameters of the node named together with the constants already
built, and \emph{none} of them from $\varepsilon,\Omega,\Omega',f,u,g$:
\begin{itemize}
\item $\delta'$, $\alpha_2$, $\lambda_2$, $R_3$, $\tau_2$, $\theta_2$,
  $\lambda_0$ and $\beta$ --- Step 1, already fixed, and fixed even before
  $c,d,M,W,\sigma$;
\item $N_1$ and $N_2$ --- the head of the application of
  \noderef{MinimumMorreyRegularity} in (2a$'$), that node binding the pair in
  its own head above the block that quantifies its data;
\item $C_1$ --- the head of the application of
  \noderef{UniformHolderDatumOnBalls} in (2a$'$), from $n,p,\tau_2,d,M$ and
  $N_2$;
\item $R_1'$ --- the head of clause (B) of
  \noderef{ComparisonGradientEstimateMin} in (2d), which is the one place
  (F4) of the thesis is consumed and the reason the case split above is made where it
  is; and then $R_1=\min\{d/16,R_1',1\}$;
\item $c_2=2^{p/2}(\omega_n+N_1^p)$ --- the formula in (2c), from $n,p$ and
  $N_1$;
\item $c_{20}$ --- the head of clause (B) of
  \noderef{ComparisonSchemeCampanatoCore} in Step 3, from $c_2$, that clause
  binding $c_{20}$ after $\lambda_2$, $R_3$ and $c_2$ and above its own data
  block;
\item $R=\min\{R_1,R_3\}$ and $\rho_0=\tfrac12R^{1+\theta_2}$ --- the two
  formulas at the end of Step 3, from $R_1$, $R_3$ and $\theta_2$;
\item $N$ --- the head of the application of \noderef{GradientDecayToC1Holder}
  in Step 4, from $n,p,\lambda_0,c_{20},\rho_0$ and $W$.
\end{itemize}
Each of those five node applications has its hypotheses split by that node's own
binder order into a \emph{head} half, discharged where the constant is produced,
and a \emph{data} half.  \emph{The head halves are what is being performed
now}, and each item below marks the point at which its data half begins.

Having fixed $N$ in this way, fix $\varepsilon,\Omega,\Omega',f,u,g$ satisfying
all the hypotheses of the statement --- (F4) of the thesis among them, consistently
with the branch.  The items (2a)--(2d) and Steps 3 and 4 below discharge, for
that data, the data halves of those same applications and read off their
conclusions.
\begin{enumerate}
\item[(2a)] \emph{The inner domain and the representative.}  Apply the
  cited H\"older estimates (C4)(a) (\noderef{CitedHolderOscillationMinimum}) with
  $\lambda,\Lambda,a,M$ --- its four guards there being $0<\lambda$,
  $\lambda\le\Lambda$, $0<a$ and $0\le M$, all hypotheses of the present
  theorem --- to obtain $\alpha_4\in(0,1]$ and $d_0>0$, which,
  after the cycle-54 binder move, are $\varepsilon$-uniform.  Then apply
  \noderef{InnerDomainRepresentative} with $n,p,\alpha_4,d_0,M,d$ and the data
  $\Omega,\Omega',u,g$.  Its hypotheses are $1\le n$ (from $2\le n$), $1\le p$
  (from $1<p$), $0<\alpha_4$ and $\alpha_4\le1$ and $0<d_0$ (all three just
  produced), $0\le M$ and $0<d$; the openness of $\Omega$; the nonemptiness and
  the boundedness of $\Omega'$; the separation
  $d\le\operatorname{dist}(x,y)$ for $x\in\Omega'$ and $y\notin\Omega$;
  $u\in W^{1,p}(\Omega)$ with weak gradient $g$ (\noderef{MemW1pOn}); and
  $|u|\le M$ almost everywhere on $\Omega$ --- every one of these except the
  three just produced being a hypothesis of the present theorem.  Its remaining
  hypothesis is the oscillation decay on every ball whose closed ball lies in
  $\Omega$, and that is exactly the clause of
  (C4) just read, at the present $\varepsilon$; its own remaining hypotheses are
  $0\le\varepsilon$ and $\varepsilon\le1$, the
  openness of $\Omega$, the growth condition (F1) of the thesis
  (\noderef{VariationalGrowthF1}), whose three fields
  \texttt{meas\_x}, \texttt{cont\_zh} and \texttt{F1} are the first three fields
  of \noderef{VariationalStructureF123} verbatim, the local minimum property
  (\noderef{IsLocalMinimumOn}) and $|u|\le M$ a.e., together with the per-ball
  guards $\overline B_{R_3}(y)\subseteq\Omega$, $0<\rho$ and $\rho\le R_3$ that
  \noderef{InnerDomainRepresentative} itself quantifies over and hands down
  unchanged.  It produces an open bounded
  $\Omega_1$ with $\overline{\Omega'}\subseteq\Omega_1$,
  $\overline{\Omega_1}\subseteq\Omega$, $d/2\le\operatorname{dist}(x,y)$ for
  $x\in\Omega'$ and $y\notin\Omega_1$, and a representative $U$ with $U=u$ a.e.\
  on $\Omega$, $U\in W^{1,p}(\Omega)$ with weak gradient $g$
  (\noderef{MemW1pOn}) and $U$ continuous on $\Omega_1$.  \emph{Three of its
  eleven conclusions are not used, and it is worth saying which.}  The
  $C^{\alpha_4}$ membership on balls of radius $d/16$ centred in $\Omega_1$
  (\noderef{MemHolderOn}) is \emph{deliberately} declined, because its exponent
  $\alpha_4$ comes from (C4) and so depends on $M$, whereas $\lambda_0$ is bound
  before $M$ --- (2a$'$) buys the same kind of datum at the fixed exponent
  $\tfrac12$ instead, from part (ii) of Theorem~3 through (C5).  The remaining two,
  $d/2\le\operatorname{dist}(x,y)$ for $x\in\Omega_1$, $y\notin\Omega$ and
  $\overline B_{d/2}(x)\subseteq\Omega$ for $x\in\Omega_1$, are simply not
  needed: this proof separates $\Omega'$ from $\partial\Omega$ directly from the
  theorem's own hypothesis, at the scale $d$ rather than $d/2$.  Note for later
  that $B_{d/2}(x)\subseteq\Omega_1$ for every
  $x\in\Omega'$: a point $y$ with $|y-x|<d/2$ cannot lie outside $\Omega_1$, by
  the separation clause $d/2\le\operatorname{dist}(x,y)$ for $x\in\Omega'$,
  $y\notin\Omega_1$ just listed.
\item[(2a$'$)] \emph{A H\"older datum at the fixed exponent $\tfrac12$, with a
  constant uniform in the centre.}  \emph{This item carries two applications,
  and each splits as Step 2 describes.}  Their head halves --- the head of
  \noderef{MinimumMorreyRegularity}, which produces $N_1$ and $N_2$, and the
  head of \noderef{UniformHolderDatumOnBalls}, which produces $C_1$ --- use only
  $n,p,\tau_2,\lambda,\Lambda,a,\nu_0,\nu_1,d,M,W,\sigma$ and the head's cited
  hypotheses, and so belong to the head half of Step 2, where they were
  performed before the data were fixed; the two sentences beginning
  \emph{``None of these head parameters''} and \emph{``None of those head
  parameters''} below mark where each one ends.  What is read at the data now
  fixed, and at a centre $y_0$ chosen inside that reading, are the two data
  halves.

  Fix $y_0\in\Omega'$ and put
  $B^{y_0}=B_{d/4}(y_0)$.  Then $B^{y_0}$ is
  open, nonempty and bounded; $\overline{B^{y_0}}=\overline B_{d/4}(y_0)
  \subseteq\Omega$, because $|y-y_0|\ge d$ for every $y\notin\Omega$ by the
  separation hypothesis; and
  $\operatorname{dist}(z,y)\ge\operatorname{dist}(y_0,y)-d/4\ge
  d-d/4\ge d/4$ for $z\in B^{y_0}$ and $y\notin\Omega$.  \emph{No covering of
  $\overline{\Omega'}$ is taken: the ball is attached to the centre under
  consideration, and the constants below are shown to be independent of it.}
  Apply
  \noderef{MinimumMorreyRegularity} (Theorem~3) with $n,p$, the exponent
  $\tau_2$ of (1d), the parameters
  $\lambda,\Lambda,a,\nu_0,\nu_1$, the separation constant $d/4$, the bounds
  $M$ and $W$ and the modulus $\sigma$.  Its guards before the two norms are
  $2\le n$, $1<p$, $0<\tau_2$, $\tau_2<n$, $0<\lambda$, $\lambda\le\Lambda$,
  $0<a$, $0<\nu_0$, $\nu_0<\nu_1$, $0<d/4$, $0\le M$ and $0\le W$, followed by
  \noderef{CitedSobolevPoincareBall} and the universal (C2)
  (\noderef{CitedDirichletW1p}), (C1) (\noderef{CitedRegularityEpsPos}) and (C3)
  (\noderef{CitedDirichletHolder}) clauses, (C1) alone also quantifying over
  every open set; all are hypotheses of the present theorem except the two on
  $\tau_2$, which (1d) supplies, and $0<d/4$, which follows from $0<d$.
  \emph{None of these head parameters involves $y_0$, so the pair $N_1,N_2$ the
  node produces is one pair serving every centre; that is what (2c) and the
  \textup{(C5)} step below use, and it is the whole reason this route fixes the
  binder order.}
  Now with the data
  $\varepsilon,\Omega,B^{y_0},f,u,g$: its remaining hypotheses are $0\le\varepsilon$
  and $\varepsilon\le1$, the openness of
  $\Omega$, the five displayed properties of $B^{y_0}$ --- that it is open, that it
  is nonempty, that it is bounded, that $\overline{B^{y_0}}\subseteq\Omega$, and the
  separation of $B^{y_0}$ from the complement of $\Omega$ by $d/4$ --- the fact that $B^{y_0}$ is a bounded Lipschitz domain
  (\noderef{IsBoundedLipschitzDomain}), which \noderef{BallIsLipschitzDomain}
  supplies for every open ball of positive radius --- its hypotheses there are
  $1\le n$, from $2\le n$, and $0<d/4$, from $0<d$ --- and which Theorem~3 now
  requires of the set in its $\Omega'$ slot, the instance
  of (C5) at the
  exponent $n/(n-1)$ on the ball $B^{y_0}$, which is the present theorem's
  universal (C5)-on-balls hypothesis (\noderef{CitedCampanatoEmbedding}) and
  whose own Lipschitz antecedent at that ball is the same application of
  \noderef{BallIsLipschitzDomain}, the
  cited H\"older estimates (C4) (\noderef{CitedHolderOscillationMinimum}),
  (F1) of the thesis--(F3) of the thesis (\noderef{VariationalStructureF123}),
  $u\in W^{1,p}(\Omega)$ with weak gradient $g$ (\noderef{MemW1pOn}),
  $|u|\le M$ a.e., the $W^{1,p}$ bound and the local minimum property
  (\noderef{IsLocalMinimumOn}) --- the last six being hypotheses of the present
  theorem, and the two on $\varepsilon$ as well.

  \emph{What this proof takes from Theorem~3 is parts \textup{(i)} and
  \textup{(ii)}, and \emph{not} part \textup{(iii)}, even though part
  \textup{(iii)} is the H\"older statement.}  Part (iii) reads ``for each
  $0\le\eta<1$ there is $w$ with $u=w$ a.e.\ on $\Omega'$ and
  $w\in C^{\eta}(\overline{\Omega'})$ (\noderef{MemHolderOn})'', and the
  constant of that $\noderef{MemHolderOn}$ is existentially bound \emph{inside}
  the clause, after $\Omega'$ and after $u$ --- unlike parts (i) and (ii), whose
  norms $N_1,N_2$ Theorem~3 binds in its head.  The present theorem binds its
  $N$ before $\varepsilon,\Omega,\Omega',f,u,g$, and $N$ depends on $C_1$
  through $R_1$ and $R$; so a $C_1$ assembled from part (iii)'s constants would
  be bound after the data and the quantifier order of the statement would fail
  (pm-0007).  Part (ii) carries an early-bound norm and the paper's own route
  from (ii) to a H\"older bound is (4.10) of the thesis, i.e.\ the cited result (C5);
  that is the route taken below.

  \emph{Part \textup{(ii)}.}  The same application gives
  \[
  \|u\|_{\mathcal L^{q,\,\lambda^\natural}(B^{y_0})}\le N_2 ,
  \qquad
  q=\frac n{n-1},\quad \lambda^\natural=n+\Bigl(1-\frac{\tau_2}p\Bigr)q
  \tag{4.15b}
  \]
  (\noderef{CampanatoBoundOn}), with the $N_2$ of this application's head.

  \emph{Part \textup{(i)}, recorded here because \textup{(2c)} consumes it.}
  The same application gives
  \[
  \int_{B^{y_0}\cap B_\rho(x)}|\nabla u|^p\,dx\le N_1^{\,p}\rho^{\,n-\tau_2}
  \qquad\bigl(x\in B^{y_0},\ 0<\rho<\operatorname{diam}B^{y_0}=d/2\bigr),
  \tag{4.15a}
  \]
  with $N_1$ the first norm of this application
  (\noderef{MorreyBoundOn}); the guards on $x$ and $\rho$ are the ones that
  node's definition carries.  Only the second clause of
  \noderef{MorreyBoundOn} is read here; its
  first clause, $|\nabla u|\in L^p(B^{y_0})$, is not needed for (2c).

  \emph{The H\"older datum at the fixed exponent $\tfrac12$, from
  \noderef{UniformHolderDatumOnBalls}.}  The centre $y_0\in\Omega'$ above was
  arbitrary, so (4.15b) holds at every $y_0\in\Omega'$ with the one
  $N_2$ of Theorem~3's head.  Apply
  \noderef{UniformHolderDatumOnBalls} with the parameters $n,p$, the exponent
  $\tau_2$ of (1d), the separation constant $d$, the bound $M$ and that $N_2$.
  Its head guards are $2\le n$, $1<p$, $0<\tau_2$, $\tau_2\le p/2$, $0<d$,
  $0\le M$ and $0\le N_2$: $2\le n$, $1<p$, $0<d$ and $0\le M$ are hypotheses of
  the present theorem; the two on $\tau_2$ are (1d), which fixes
  $\tau_2=\min\{\tfrac12p\alpha_2\theta_2,\tfrac n2,\tfrac p2\}$ and so gives
  both $\tau_2>0$ and $\tau_2\le p/2$; and $0\le N_2$ is part of Theorem~3's
  head.  Its remaining head hypothesis is the cited Campanato embedding (C5) at
  the exponent $q=n/(n-1)$ on the \emph{single} ball $B_{d/4}(0)$
  (\noderef{CitedCampanatoEmbedding}), which is the present theorem's universal
  (C5)-on-balls hypothesis read at the centre $0$ and the radius $d/4$, whose
  guard $0<d/4$ follows from $0<d$.

  None of those head parameters involves $\varepsilon,\Omega,\Omega',f,u,g$ or
  the centre $y_0$, so the constant $C_1\ge0$ that node produces is bound before
  the data, which is what the present theorem needs: $N$ depends on $C_1$
  through $R_1$ and $R$, and $N$ is bound before
  $\varepsilon,\Omega,\Omega',f,u,g$.  By that node $C_1$ is a function of
  $n,p,\tau_2,d,M$ and $N_2$ alone, hence --- by (1d) for $\tau_2$ and by
  Theorem~3's head for $N_2$ --- of
  $n,p,\lambda,\Lambda,a,\nu_0,\nu_1,\delta,c,d,M,W$ and $\sigma$.

  Now supply that node's data hypotheses, with the present $\Omega$ and
  $\Omega'$, with $u$ in the role of its $u$ and with the $U$ of (2a) in the
  role of its $U$:
  \begin{itemize}
  \item $\Omega$ is open, and $d\le\operatorname{dist}(x,y)$ for every
    $x\in\Omega'$ and every $y\notin\Omega$: both are hypotheses of the present
    theorem;
  \item $|u|\le M$ almost everywhere on $\Omega$: a hypothesis of the present
    theorem;
  \item $U=u$ almost everywhere on $\Omega$: this is (2a);
  \item $U$ is continuous on $B_{d/4}(y_0)$ for every $y_0\in\Omega'$: by (2a),
    $U$ is continuous on $\Omega_1$, and
    $B_{d/4}(y_0)\subseteq B_{d/2}(y_0)\subseteq\Omega_1$ for every
    $y_0\in\Omega'$, the first inclusion because $d/4\le d/2$ and the second
    being the one noted at the end of (2a); and continuity on a set restricts
    to every subset of it;
  \item $\|u\|_{\mathcal L^{q,\lambda^\natural}(B_{d/4}(y_0))}\le N_2$ for every
    $y_0\in\Omega'$ (\noderef{CampanatoBoundOn}), at the exponent
    $\lambda^\natural=n+\bigl(1-\tfrac{\tau_2}p\bigr)q$: this is
    (4.15b), which is exactly the exponent and the constant that
    display carries.
  \end{itemize}
  Its conclusion is that for every $x_0\in\Omega'$ and every $R$ with
  $0<R\le d/16$ one has $U\in C^{1/2}(\overline B_R(x_0))$
  (\noderef{MemHolderOn}) and
  \[
  |U(x)-U(y)|\le C_1\operatorname{dist}(x,y)^{1/2}
  \qquad\bigl(x,y\in\overline B_R(x_0)\bigr),
  \]
  with the \emph{single} constant $C_1$ for every such centre and radius.  This
  is both the H\"older datum clause (B) of
  \noderef{ComparisonSchemeCampanatoCore} asks for at $\alpha_H=\tfrac12$ and
  the two-point datum clause (B) of \noderef{ComparisonGradientEstimateMin}
  asks for at $\alpha=\tfrac12$ with its constant $C_1$, which that clause binds
  before its radius and therefore before the centre --- and which, by the
  paragraph above, $C_1$ is eligible for.
\item[(2b)] \emph{The frozen family.}  Put
  $\mathcal A(x,z)=\nabla_hf(x,z,\cdot)$
  (\noderef{FrozenGradientOperator}).  For every $x_0\in\Omega'\subseteq\Omega$
  --- the inclusion because $\Omega'\subseteq\overline{\Omega'}\subseteq\Omega$
  --- apply \noderef{FrozenOperatorStructureQ12} with
  $n,p,\varepsilon,\lambda,\Lambda,a,\nu_0,\nu_1$, the open set $\Omega$, the
  density $f$, the modulus $\sigma$, the point $x_1=x_0$ and the value
  $z_1=U(x_0)$.  Its hypotheses are the present
  theorem's $2\le n$, $1<p$, $0\le\varepsilon$ and $\varepsilon\le1$,
  $0<\lambda$, $\lambda\le\Lambda$,
  $0<a$, $0<\nu_0$ and $\nu_0<\nu_1$, the openness of $\Omega$, $x_0\in\Omega$
  and (F1) of the thesis--(F3) of the thesis (\noderef{VariationalStructureF123}); its
  value $z_1$ is universally quantified, so $U(x_0)$ is an admissible choice.

  \emph{Its conclusion is the whole bundle, not the two inequalities, and that
  is what clause \textup{(B)} needs.}  The conclusion is
  $\mathrm{EllipticStructureQ12}\;n\;p\;\varepsilon\;\nu_0\;\nu_1\;\mathrm{univ}$
  for the normalised field
  $\mathcal A(x_0,U(x_0))-\mathcal A(x_0,U(x_0))(0)$
  (\noderef{EllipticStructureQ12}) --- all six of that predicate's clauses, the
  two Carath\'eodory clauses and \texttt{hcont} and \texttt{hdiff} as well as
  (Q1) of the thesis and (Q2) of the thesis, delivered in one piece, with the structure constants
  $\gamma_0=\nu_0$ and $\gamma_1=\nu_1$ undegraded and the normalisation
  $A\mapsto A-A(0)$ of thesis l.~542 already performed inside the node.  Clause
  (B) of \noderef{ComparisonSchemeCampanatoCore} asks for exactly that
  predicate, so no clause of it has to be assembled here.  \emph{This is why
  the present proof needs no promotion of the a.e.-$x$ clauses
  \textup{(Q1)}--\textup{(Q2)} to a single point (pm-0034): on the variational
  side those two inequalities arrive as a \emph{conclusion} built from
  (F1) of the thesis--(F3) of the thesis, not as a hypothesis read at $x_0$.}  The guard
  $\gamma_1\le r\gamma_0$ of clause (B) holds with equality by (1a).
\item[(2c)] \emph{The Morrey bound (4.15) of the thesis.}  The paper's step here is
  ``given $\tau>0$, by Theorem~Theorem~3 of the thesis there exists a constant $c_1$ such
  that $\Lambda(x_0,R)\le c_1R^{-\tau}$ for $R\le R_1$'' (l.~1371--1374), and
  Theorem~3 is applied to $\Omega'$.  \emph{Read against
  \noderef{MinimumMorreyRegularity}'s statement, that application does not
  deliver what clause \textup{(B)} asks for, and the gap is in the domain of
  integration, not in the exponent.}  Part (i) of that node is a
  \noderef{MorreyBoundOn} assertion about its own $\Omega'$ argument, and
  \noderef{MorreyBoundOn} bounds
  $\int_{\Omega'\cap B_\rho(x)}|\nabla u|^p$ --- the integral over the
  \emph{intersection}.  Clause (B) of \noderef{ComparisonSchemeCampanatoCore}
  asks for $R^{-n}\int_{B_R(x_0)}(1+|\nabla u|^2)^{p/2}\le c_2R^{-\tau_2}$, the
  integral over the \emph{whole} ball; and for $x_0\in\Omega'$ the ball
  $B_R(x_0)$ need not be contained in $\Omega'$, which is a bounded Lipschitz
  domain and not a neighbourhood of its own boundary points.  So the bound is
  taken where the containment \emph{is} available, namely at the ball of
  (2a$'$) centred at the same point.

  Let $x_0\in\Omega'$ and $0<R\le d/16$, and take
  the ball $B^{x_0}=B_{d/4}(x_0)$ of (2a$'$) at $y_0=x_0$; then
  $\overline B_R(x_0)\subseteq B^{x_0}$ since $R\le d/16<d/4$, so
  $x_0\in B^{x_0}$, $B^{x_0}\cap B_R(x_0)=B_R(x_0)$ and
  $0<R\le d/16<d/2=\operatorname{diam}B^{x_0}$.  Hence (4.15a)'s two guards, on its
  centre and on its radius, are met at $x=x_0$ and $\rho=R$, and it reads
  \[
  \int_{B_R(x_0)}|\nabla u|^p\,dx\le N_1^{\,p}R^{\,n-\tau_2} ,
  \]
  with the single centre-independent $N_1$ of (2a$'$).  Therefore, by
  $(1+t^2)^{p/2}\le2^{p/2}(1+t^p)$ for $t\ge0$ and $|B_R|=\omega_nR^n$
  (Preamble: the volume of a ball), $\omega_n$ being the volume of the unit
  ball,
  \[
  \Lambda(x_0,R)=R^{-n}\!\int_{B_R(x_0)}(1+|\nabla u|^2)^{\frac p2}dx
  \le2^{\frac p2}\bigl(\omega_n+N_1^pR^{-\tau_2}\bigr)\le c_2R^{-\tau_2}
  \]
  for $x_0\in\Omega'$ and $0<R\le\min\{1,d/16\}$, with
  $c_2=2^{\frac p2}(\omega_n+N_1^p)$; the last step uses $R^{-\tau_2}\ge1$ for
  $R\le1$, which is where $R\le1$ enters.  This
  is clause (B)'s fifteenth hypothesis, and $c_2$ is passed as clause (B)'s
  $c_2$, whose guard $0<c_2$ holds because $\omega_n>0$ and $N_1^p\ge0$.  Clause
  (B) binds $c_2$ \emph{after} $\lambda_2$ and $R_3$, so the $d$, $M$, $W$ and
  $\sigma$ inside $N_1$ reach only the $c_{20}$ it then produces and never
  $\lambda_0$.  ($N_1$ does
  carry $\sigma$, because Theorem~3 binds $\sigma$ before its $N_1$; this is why
  the present statement binds $\sigma$ before $N$ and after $\lambda_0$.)

  \emph{The per-instance \textup{(C5)} hypothesis on $\Omega'$ at the exponent
  $n/(n-1)$, and the second application of Theorem~3 it carries.}  Apply
  \noderef{MinimumMorreyRegularity} once more, with the same head parameters
  except the separation constant $d$ in place of $d/4$, and with the data
  $\varepsilon,\Omega,\Omega',f,u,g$: its inner hypotheses are the same list as
  in (2a$'$) read at the present $\Omega'$, every one of them a hypothesis of
  the present theorem --- the openness, nonemptiness and boundedness of
  $\Omega'$, $\overline{\Omega'}\subseteq\Omega$ and the separation
  $d\le d(\Omega',\partial\Omega)$ among them.  Two of those hypotheses are worth
  naming separately, because they are two different hypotheses of the present
  theorem and not one.  First, the Lipschitz hypothesis on $\Omega'$ that
  Theorem~3 now carries is the present theorem's own: $\Omega'$ is a bounded
  Lipschitz domain (\noderef{IsBoundedLipschitzDomain}).  Second, the instance
  of (C5) at the exponent $n/(n-1)$ on $\Omega'$ that Theorem~3 also carries is
  the present theorem's \emph{per-instance} one on $\Omega'$ at that exponent
  (\noderef{CitedCampanatoEmbedding}), as distinct from the per-instance one at
  the exponent $p$ consumed in Step 4 and from the universal one on balls
  consumed in (2a$'$); this second application of Theorem~3 is where that
  hypothesis is consumed.  Its part (i) is the paper's (4.15) of the thesis in the
  form Theorem~3 states it,
  $\int_{\Omega'\cap B_\rho(x)}|\nabla u|^p\le N_1'^{\,p}\rho^{\,n-\tau_2}$
  (\noderef{MorreyBoundOn}).  It is recorded because it is the paper's own
  citation at this point; the bound clause (B) consumes is the one displayed
  above, for the reason given there.
\item[(2d)] \emph{The smallness (4.17) of the thesis.}  Read clause (B) of
  \noderef{ComparisonGradientEstimateMin} --- the instance of (1b), at
  $\alpha=\tfrac12$ and the present $\delta$ --- at the present $M$, $d$ and
  $\sigma$ and at the constants $c_\sigma=c$ of (F4) of the thesis and
  $C_1$ of (2a$'$).  The guards this clause binds before $R_1'$ are $0<\alpha$
  and $\alpha\le1$, which hold at $\alpha=\tfrac12$, $0<\delta$, a hypothesis of
  the present theorem, then $0\le M$ and $0<d$, likewise hypotheses, then
  $0<c_\sigma$, which is the present theorem's $0<c$, and $0\le C_1$, which
  (2a$'$) records, and finally $\sigma(t)\le c\,t^{\delta}$ for $t\ge0$, which is
  (F4) of the thesis --- decided, in this branch, at the head of Step 2, which is the
  one place that condition is consumed.  \emph{Those are all of the hypotheses
  the clause binds above $R_1'$}, so the radius it produces is a function of
  $n,p,\lambda,\Lambda,a,\nu_0,\nu_1,\delta$ and $M,d,\sigma,c,C_1$ alone, with
  no datum among them; let $R_1'>0$ be that radius and put
  \[
  R_1=\min\{d/16,\,R_1',\,1\}>0 ,
  \]
  both fixed here, before $\varepsilon,\Omega,\Omega',f,u,g$.  Then
  $R_1\le R_1'$, which is what the display below needs, and \emph{$R_1\le1$,
  which is what} (2c) \emph{needs} --- without the third term in the minimum,
  $d/16$ and $R_1'$ could both exceed $1$ and the Morrey bound of (2c), which
  holds only for $R\le1$, would not cover every $R\le R_1$.

  \emph{Now with the data.}  \emph{Apply the clause's inner block with $U$, not
  $u$, in the role of its
  $W^{1,p}$ function}, which is legitimate and is what makes the bundles match:
  $U=u$ a.e.\ on $\Omega$ by (2a), so $|U|\le M$ a.e.\ on $\Omega$, and
  $U\in W^{1,p}(\Omega)$ with weak gradient $g$ (\noderef{MemW1pOn}) and $U$ is
  a local minimum for $F$ in $\Omega$ (\noderef{IsLocalMinimumOn}) by
  \noderef{AEModificationInvariance}, applied with $q=p$, the open set $\Omega$,
  the pair $(u,U)$ and the field $g$ --- its hypotheses are $1\le n$ (from
  $2\le n$), $1\le p$ (from $1<p$), the openness of $\Omega$ and the
  almost-everywhere equality $u=U$ on $\Omega$, which is (2a)'s $U=u$ read with
  the two sides exchanged --- and of its five conclusions the two used are the
  second, transporting $\mathrm{MemW1pOn}$, and the fifth, transporting
  $\mathrm{IsLocalMinimumOn}$ at the present $f$.  Its first, third and fourth
  --- transporting $\mathrm{HasWeakGradientOn}$, $\mathrm{MemW1pLocOn}$ and
  $\mathrm{IsWeakSolutionOn}$ --- are not used: the first and third are weaker
  than the second for the present purpose, and the fourth belongs to the
  equation side rather than the variational one.  With that choice the comparison
  lemma's
  $\mathrm{FrozenComparisonOn}$ argument list reads
  $(n,p,x_0,R,\mathcal A(x_0,U(x_0)),U,U,g,v,V,g_v)$, which is the \emph{same}
  list clause (B) of \noderef{ComparisonSchemeCampanatoCore} quantifies over in
  its fifth hypothesis; had we applied it with $u$, the bundle's clause
  ``$U=u$ on $B_R(x_0)$'' would be a pointwise identity that only holds almost
  everywhere here, and the two hypotheses would not meet.  The inner block then
  gives, for every $x_0\in\Omega'$, every $0<R\le R_1'$ with
  $\overline B_R(x_0)\subseteq\Omega$ and every
  $\mathrm{FrozenComparisonOn}(n,p,x_0,R,\mathcal A(x_0,U(x_0)),U,U,g,v,V,g_v)$
  (\noderef{FrozenComparisonOn}),
  \[
  \int_{B_R(x_0)}|\nabla u-g_v|^p\,dx\le R^{\delta'}\!\int_{B_R(x_0)}(1+|\nabla
  u|^2)^{\frac p2}dx=R^{\delta_2+n}\Lambda(x_0,R),
  \]
  which is clause (B)'s sixteenth hypothesis.  \emph{Its own remaining
  hypotheses are \textup{not} the list of} (2a$'$), \emph{which is a different
  node's list read at the ball $B^{y_0}$; they are the following, every one of them a
  hypothesis of the present theorem except where stated.}
  $0\le\varepsilon$ and $\varepsilon\le1$; the openness of $\Omega$ and of
  $\Omega'$, the boundedness of $\Omega'$ and
  $\overline{\Omega'}\subseteq\Omega$; the separation
  $d\le\operatorname{dist}(x,y)$ for $x\in\Omega'$ and $y\notin\Omega$;
  (F1) of the thesis--(F3) of the thesis at the present parameters and modulus
  (\noderef{VariationalStructureF123}); the cited H\"older estimates (C4)
  (\noderef{CitedHolderOscillationMinimum}); $U\in W^{1,p}(\Omega)$ with weak gradient
  $g$ (\noderef{MemW1pOn}), $|U|\le M$ a.e.\ and $U$ a local minimum for $F$ in
  $\Omega$ (\noderef{IsLocalMinimumOn}) --- these three by (2a) and the transfer
  just made, since the node is being read with $U$ in the role of its $u$;
  $x_0\in\Omega'$, $0<R$, $R\le R_1'$ and
  $\overline B_R(x_0)\subseteq\Omega$, which are the per-centre guards the
  clause quantifies over and which the choice of $R_1$ below supplies; the
  frozen comparison bundle itself (\noderef{FrozenComparisonOn}), which clause
  (B) of \noderef{ComparisonSchemeCampanatoCore} also quantifies over and hands
  in; and the two-point H\"older bound for $U$ on
  $\overline B_R(x_0)$ at the exponent $\tfrac12$ and the constant $C_1$, which
  (2a$'$) supplies for every $x_0\in\Omega'$ and every $R\le d/16$.  Note that
  the head parameters $n,p,\lambda,\Lambda,a,\nu_0,\nu_1$ and the two head
  hypotheses \noderef{CitedSobolevPoincareBall} and the universal (C3)
  (\noderef{CitedDirichletHolder}) were already discharged in (1b), where this
  same instance was created; no second instance is taken here.  Finally, with
  the $R_1$ fixed above, $\overline B_R(x_0)\subseteq\Omega$ for every
  $x_0\in\Omega'$ and every $0<R\le R_1$, since $R_1\le d/16<d$ and the
  separation hypothesis gives $|y-x_0|\ge d$ for $y\notin\Omega$.
\end{enumerate}

\emph{Step 3: the decay (4.19) of the thesis.}  Read clause (B) of
\noderef{ComparisonSchemeCampanatoCore} --- the instance produced in (1d) ---
at the $c_2$ of (2c), whose guard $0<c_2$ holds there.  That clause binds
$c_{20}$ immediately after $c_2$, hence after $\lambda_2$ and $R_3$, and
\emph{above} the block that quantifies its own data, so the constant $c_{20}>0$
it produces is a function of the head parameters and $c_2$ alone and is fixed
here, before $\varepsilon,\Omega,\Omega',f,u,g$, as the order of Step 2
requires.

\emph{Now with the data} $\varepsilon,\nu_0,\nu_1,\Omega,\Omega',\mathcal
A,u,U,g,R_1$.  \emph{That clause binds sixteen hypotheses, and
listing nine of them and gesturing at the rest is what hid the two defects
repaired in \textup{(2c)} and \textup{(2d)}; collected in the node's own order
they are:}
\begin{itemize}
\item $0\le\varepsilon$ and $\varepsilon\le1$ --- hypotheses of the present
  theorem;
\item $0<\gamma_0=\nu_0$ and $0<\gamma_1=\nu_1$ --- the second from
  $0<\nu_0<\nu_1$;
\item $\gamma_1\le r\gamma_0$, i.e.\ $\nu_1\le r\nu_0$, which holds with
  equality by the definition $r=\nu_1/\nu_0$ of (1a);
\item the openness of $\Omega$ and of $\Omega'$, and
  $\overline{\Omega'}\subseteq\Omega$ --- hypotheses of the present theorem;
\item $0<R_1$, by the choice $R_1=\min\{d/16,R_1',1\}$ at the end of (2d),
  all three terms being positive;
\item $u\in W^{1,p}(\Omega)$ with weak gradient $g$ (\noderef{MemW1pOn}), a
  hypothesis of the present theorem, and $U\in W^{1,p}(\Omega)$ with the
  \emph{same} weak gradient $g$ (\noderef{MemW1pOn}), by (2a);
\item $\overline B_R(x_0)\subseteq\Omega$ for every $x_0\in\Omega'$ and every
  $0<R\le R_1$, which is the inclusion recorded at the end of (2d);
\item for every $x_0\in\Omega'$, the bundle
  (Q1) of the thesis--(Q2) of the thesis on $\R^n$ at $p,\varepsilon,\nu_0,\nu_1$ for the
  normalised frozen field
  $\mathcal A(x_0,U(x_0))-\mathcal A(x_0,U(x_0))(0)$
  (\noderef{EllipticStructureQ12}), which is (2b) --- read at $x_0$, which lies
  in $\Omega$ because $\Omega'\subseteq\overline{\Omega'}\subseteq\Omega$;
\item for every $x_0\in\Omega'$, the H\"older datum
  $U\in C^{\alpha_H}(\overline B_{R_1}(x_0))$ at $\alpha_H=\tfrac12$
  (\noderef{MemHolderOn}), which is (2a$'$) read at the single radius $R_1$,
  legitimate because $R_1\le d/16$; the constant $C_1$ there is uniform in the
  centre, which is what this per-centre clause requires, since
  \noderef{MemHolderOn} binds its constant inside and the clause is quantified
  over $x_0$ outside;
\item the Morrey hypothesis
  $R^{-n}\int_{B_R(x_0)}(1+|\nabla u|^2)^{p/2}\le c_2R^{-\tau_2}$ for every
  $x_0\in\Omega'$ and every $0<R\le R_1$, which is (2c), using $R_1\le1$ and
  $R_1\le d/16$; and
\item the smallness hypothesis, for every $x_0\in\Omega'$, every
  $0<R\le R_1$ and every frozen comparison bundle
  $\mathrm{FrozenComparisonOn}(n,p,x_0,R,\mathcal A(x_0,U(x_0)),U,U,g,v,V,g_v)$
  (\noderef{FrozenComparisonOn}), which is (2d), using $R_1\le R_1'$.
\end{itemize}
It gives, for every
$x_0\in\Omega'$, every $0<R\le\min\{R_1,R_3\}$ and every $\rho>0$ with
$2\rho\le R^{1+\theta_2}$
\[
\int_{B_\rho(x_0)}\bigl|\nabla u-(\nabla u)_{B_\rho(x_0)}\bigr|^pdx
\le c_{20}\,\rho^{\,n+\lambda_2}.
\tag{4.19}
\]
Fix
\[
R=\min\{R_1,R_3\},\qquad \rho_0=\tfrac12R^{1+\theta_2}
\]
once, so that (4.19) of the thesis holds for all $0<\rho\le\rho_0$.  \emph{Both are
built from $R_1$ of \textup{(2d)}, $R_3$ of \textup{(1d)} and $\theta_2$ of
\textup{(1d)}, so both belong to the head half and were available before the
data were fixed}, which is what Step 2's order asserts of them; they are written
here only because this is where (4.19) of the thesis is read.

\emph{Step 4: parts \textup{(i)} and \textup{(ii)}, from
\noderef{GradientDecayToC1Holder}.}  Apply that node with the parameters $n,p$,
with $\lambda_0$ in the role of its $\lambda_0$ --- recall $\lambda_0=\lambda_2$
by (1e) --- with the $c_{20}$ of Step 3 in the role of its $c_{20}$, with the
$\rho_0=\tfrac12R^{1+\theta_2}$ of Step 3 in the role of its $\rho_0$, and with
$W$ in the role of its $W$.  \emph{Every one of those six head parameters was
fixed in the head half, before $\varepsilon,\Omega,\Omega',f,u,g$}, and that
node binds its constant in its own head, above the block quantifying its data;
so the $N$ it produces is fixed there too, which is exactly what the statement's
$\exists N$ above the data demands.  Its head guards hold:
\begin{itemize}
\item $2\le n$ and $1<p$, and $0\le W$, are hypotheses of the present theorem;
\item $0<\lambda_0$ and $\lambda_0<p$ are (1e), where the clause of (1d) that
  produces $\lambda_2$ asserts $0<\lambda_2<p$;
\item $0<c_{20}$ is recorded in Step 3, that clause producing $c_{20}$ with
  $c_{20}>0$;
\item $0<\rho_0$, because $R_1>0$ by (2d) and $R_3>0$ by (1d), so $R>0$, and
  $R^{1+\theta_2}>0$ for a positive base and any exponent; and $\rho_0\le1$,
  because $R\le R_3\le1$ by (1d) and $1+\theta_2>0$ give $R^{1+\theta_2}\le1$,
  whence $\rho_0\le\tfrac12\le1$.
\end{itemize}
Let $N\ge0$ be the constant it produces.

\emph{The dependence of $N$, against Theorem~Theorem~4 of the thesis's own list.}  By that
node $N$ is a function of $n,p,\lambda_0,c_{20},\rho_0$ and $W$ alone.  Of these
$\lambda_0=\lambda_2$ and $\theta_2$ are functions of the quantities (1e) names;
$\rho_0$ is built from $R=\min\{R_1,R_3\}$ and $\theta_2$, with $R_3$ from (1d),
again on that list, and $R_1$ from (2d), which carries $d$, $M$, $\sigma$, the
constant $c$ of (F4) of the thesis and the $C_1$ of (2a$'$) --- clause (B) of
\noderef{ComparisonGradientEstimateMin} binds its $R_1'$ after all five --- and
$C_1$ is, by (2a$'$), a function of $n,p,\tau_2,d,M$ and $N_2$, with $N_2$ bound
in Theorem~3's head; and $c_{20}$ carries $c_2$, hence $N_1$, hence $d$, $M$,
$W$ and $\sigma$ by (2c).  \emph{Every one of these is bound before
$\varepsilon,\Omega,\Omega',f,u,g$, which is what the statement's quantifier
order requires of $N$.}  So $N$ depends on
$n,p,\lambda,\Lambda,a,\nu_0,\nu_1,\delta$ together with $c,d,M,W$ and
$\sigma$ --- precisely what Theorem~Theorem~4 of the thesis permits at l.~1354--1355, where
the norm may depend, ``in addition to the previous constants'', on
$d(\Omega',\partial\Omega)$, $\|u\|_{L^\infty(\Omega)}$ and
$\|u\|_{W^{1,p}(\Omega)}$.  \emph{The appearance of $M$ in $N$ is licensed, and
licensed only for $N$}: $\lambda_0$ and $\beta$ are bound before $M$ by (1e),
which is the other half of the paper's remark and the reason (1b) and (1c) take
the trouble they do.

\emph{The data hypotheses of that node.}  Supply the present $\Omega,\Omega',u$
and $g$:
\begin{itemize}
\item $\Omega$ and $\Omega'$ are open and $\overline{\Omega'}\subseteq\Omega$:
  hypotheses of the present theorem;
\item $\Omega'$ is a bounded Lipschitz domain
  (\noderef{IsBoundedLipschitzDomain}): a hypothesis of the present theorem;
\item the cited Campanato embedding (C5) at the exponent $p$ on $\Omega'$
  (\noderef{CitedCampanatoEmbedding}): a hypothesis of the present theorem ---
  this is the per-instance one at the exponent $p$, not the one at $n/(n-1)$ and
  not the universal one on balls;
\item $u\in W^{1,p}(\Omega)$ with weak gradient $g$ (\noderef{MemW1pOn}): a
  hypothesis of the present theorem;
\item $\int_\Omega|\nabla u|^p\,dx\le W^p$: the present theorem assumes
  $\int_\Omega|u|^p\,dx+\int_\Omega|\nabla u|^p\,dx\le W^p$, and
  $\int_\Omega|u|^p\,dx\ge0$, being the integral of a nonnegative function, so
  dropping it only decreases the left-hand side;
\item the containment $B_\rho(x_0)\subseteq\Omega$ for every $x_0\in\Omega'$ and
  every $0<\rho\le\rho_0$.  This is the separation hypothesis of the present
  theorem, read at the radius $\rho_0$ of Step 3.  Indeed (2d) records
  $\overline B_R(x_0)\subseteq\Omega$ for every $x_0\in\Omega'$ and every
  $0<R\le R_1$, with $R_1=\min\{d/16,R_1',1\}$; and
  $\rho_0=\tfrac12R^{1+\theta_2}$ with $R=\min\{R_1,R_3\}$, so $R\le R_1\le1$
  and $\theta_2>0$ give $R^{1+\theta_2}\le R\le R_1$ and hence
  $\rho_0\le\tfrac12R_1<R_1$.  So for $0<\rho\le\rho_0$ we have $\rho\le R_1$
  and therefore
  $B_\rho(x_0)\subseteq\overline B_\rho(x_0)\subseteq\Omega$, an open ball being
  contained in the closed ball of the same centre and radius;
\item the decay hypothesis, namely
  $\int_{B_\rho(x_0)}\bigl|\nabla u-(\nabla u)_{B_\rho(x_0)}\bigr|^pdx\le
  c_{20}\rho^{\,n+\lambda_0}$ for every $x_0\in\Omega'$ and every
  $0<\rho\le\rho_0$: this is (4.19) of the thesis at the fixed
  $R=\min\{R_1,R_3\}$, exactly as recorded at the end of Step 3, since
  $\lambda_0=\lambda_2$ and $\rho_0=\tfrac12R^{1+\theta_2}$.
\end{itemize}

Its conclusion is the conjunction of the two assertions of the present theorem:
\begin{itemize}
\item $\|D_iu\|_{\mathcal L^{p,\,n+\lambda_0}(\Omega')}\le N$ for every
  $i=1,\dots,n$ (\noderef{CampanatoBoundOn}), which is (i); and
\item functions $w$ and $g_w$ with $u=w$ almost everywhere on $\Omega'$,
  $g=g_w$ almost everywhere on $\Omega'$, and
  $\mathrm{MemC1HolderOn}\;n\;\beta\;\Omega'\;\overline{\Omega'}\;w\;g_w$ at
  $\beta=\lambda_0/p$ (\noderef{MemC1HolderOn}) --- that is, $w$ differentiable
  at every point of $\Omega'$ with $\nabla w=g_w$ there, and $g_w$ continuous
  and $\beta$-H\"older on $\overline{\Omega'}$ --- which is (ii), with
  $\beta=\lambda_0/p$ as the statement asserts and $0<\beta\le1$ by (1e).
\end{itemize}
This completes the proof.
\end{proof}
